\chapter{Getting Started}\label{ch:start}

This chapter takes you from an empty directory to running Dylan in both
implementations. It lists what you need to install, shows how to get the
source, walks through a first session with the interactive
implementation, and compiles a first program with \cmd{dylanc}. By the
end of it you will have everything you need for the tutorial in
Chapter~\ref{ch:tutorial}.

\section{Requirements}\label{sec:start-requirements}

\index{requirements}
Both implementations are written in Common Lisp and need \textbf{SBCL},
Steel Bank Common Lisp. The system was developed on SBCL~2.6.3, and any
recent release should work. The interactive implementation needs nothing
else but \cmd{make} to build it, and even that can be done without.

The compiler to C needs two more things, and can use a third:

\begin{itemize}
\item \textbf{A C compiler that supports the \code{musttail}
  attribute}, which guarantees that a call in tail position is compiled as
  a jump. That means \textbf{GCC 15 or later} or \textbf{Clang 13 or
  later}. Any other C11 compiler can be used with the \cmd{-{}-trampoline}
  option, which implements tail calls with trampolines instead, at some
  cost in speed. \cmd{dylanc} uses the compiler named by the \code{CC}
  environment variable, or \cmd{cc} if \code{CC} is not set; the
  \cmd{-{}-cc} option names one explicitly.
\item \textbf{\cmd{make}} and a POSIX shell. \cmd{bin/dylanc} is a shell
  script, and it builds the C runtime library with \cmd{make} the first
  time it is needed.
\item Optionally, \textbf{the Boehm-Demers-Weiser garbage collector},
  \file{libgc}, \emph{built with thread support}, together with its
  header files, for programs built with \cmd{-{}-gc boehm}
  (Section~\ref{sec:dylanc-gc}); the runtime then defines
  \code{GC\_THREADS} and links with \cmd{-lgc -lm -lpthread}. The
  default collector, the Memory Pool System, is part of the source tree
  and needs nothing installed.
\end{itemize}

\paragraph{Checking what you have.} These commands report the versions
installed:

\begin{shell}
$ sbcl --version
SBCL 2.6.3
$ gcc --version | head -1
gcc (GCC) 16.2.1 20260819 (Red Hat 16.2.1-2)
$ clang --version | head -1
clang version 22.1.8 (Fedora 22.1.8-4.fc44)
\end{shell}

Your version numbers will differ; what matters is that SBCL runs and that
one of the C compilers is new enough (or that you plan to use
\cmd{-{}-trampoline}). For the Boehm collector, the header \file{gc.h}
must be present too (\cmd{ls /usr/include/gc.h}). On Fedora and related
systems its development package is called \file{gc-devel}; on Debian and
Ubuntu it is \file{libgc-dev}. The packaged collectors on these systems
are normally built with thread support. If you build the collector
yourself, configure it with threads enabled.

\begin{notebox}
The C compiler is needed only to \emph{build} programs. An executable made
by \cmd{dylanc} needs only the C library at run time (and the shared
\file{libgc}, if built with \cmd{-{}-gc boehm}); it does not need SBCL.
\end{notebox}

\section{Getting the source}\label{sec:start-source}

The project lives at \url{https://github.com/blakemcbride/Dylan}. Clone
it with Git:

\begin{shell}
$ git clone https://github.com/blakemcbride/Dylan.git
$ cd Dylan
\end{shell}

There is nothing to configure and nothing to install. Every command in
this book is run from the top directory of the clone unless the text says
otherwise. The directories you will meet first are these:

\begin{center}
\begin{tabular}{@{}ll@{}}
\toprule
\file{src/} & the interactive implementation, in Common Lisp \\
\file{Makefile} & builds \file{bin/dylan} \\
\file{bin/dylan} & the interactive implementation, once built \\
\file{bin/dylanc} & the command that runs the compiler to C \\
\file{compiler/} & the compiler to C, in Common Lisp \\
\file{runtime/} & the compiler's runtime: \file{runtime/c} in C,
  \file{runtime/dylan} in Dylan \\
\file{manual/Dylan-1992/} & the annotated 1992 Dylan manual (the DIRM) \\
\file{dylan-examples.dyl} & the examples from the 1992 manual \\
\file{tests/}, \file{bench/} & test suites and benchmarks \\
\bottomrule
\end{tabular}
\end{center}

You can run \cmd{dylan} and \cmd{dylanc} from anywhere by putting
\file{bin} on your \code{PATH}, or by making symbolic links to them in a
directory that is already on it. \cmd{dylan} is self-contained, and the
\cmd{dylanc} script finds the rest of the system relative to its own real
location, so a link works for both.

\section{The interactive implementation}\label{sec:start-rep}

\index{read-eval-print loop}\index{dylan@\cmd{dylan}}
Build the interactive implementation once, in the top directory of the
clone:

\begin{shell}
$ make
sbcl --script src/build-dylan.lisp
;Writing /home/you/Dylan/bin/dylan
\end{shell}

\noindent
This loads the whole system into SBCL and saves it as one executable,
\file{bin/dylan}. (On Windows, without \cmd{make}, run
\cmd{sbcl -{}-script src/build-dylan.lisp} yourself; it writes
\file{bin\textbackslash dylan.exe}.) Run \cmd{make} again whenever
\file{src/} changes; it does nothing when \file{bin/dylan} is up to date.

\index{fasl files}
The first build compiles every Lisp file of the implementation and
writes the compiled files (\emph{fasl} files) into a directory under
\file{src/fasl/} named for the Lisp and its version, such as
\file{src/fasl/SBCL-2-6-3/}. This takes a little while. Later builds
recompile only a changed source file and every file loaded after it,
since later files may depend on its macros; you never need to delete the
fasl directory by hand. Because the fasl files are kept apart for each
Lisp version, upgrading SBCL simply makes a new directory.

\paragraph{Starting the loop.} Run \cmd{dylan}. It starts a
read-eval-print loop that reads Dylan rather than Lisp. Its prompt is a
question mark:

\begin{repl}
$ bin/dylan

Entering Dylan read-eval-print-loop.
Leave by typing "(exit)"

? (define greeting "Hello")

;Value: greeting

? (concatenate greeting ", world")

;Value: "Hello, world"
\end{repl}

Type a form and press Enter. The loop reads the form, compiles it into
Lisp, runs it, and prints its value after \code{;Value:}. A form may
span several lines; the loop keeps reading until the parentheses balance.
A definition such as \code{define} answers the name it defined. A form
that returns several values shows each one on its own line, numbered:

\begin{repl}
? (truncate/ 17 5)

;Value[1]: 3
;Value[2]: 2
\end{repl}

\begin{notebox}
In a real terminal the loop prints an extra blank line after each form you
type. The transcripts in this book leave it out.
\end{notebox}

\paragraph{Leaving the loop.} Type \code{dylan:done} or \code{(exit)},
or end the input with Ctrl-D, to leave \cmd{dylan}.

\paragraph{Loading it into SBCL yourself.} \index{dylan:rep}
\index{load-dylan.lisp@\file{load-dylan.lisp}}
Instead of \cmd{dylan}, you can start a plain SBCL in the top directory
of the clone, load \file{src/load-dylan.lisp}, and call the function
\code{dylan:rep}, which starts the same loop. This needs no build step,
since the load compiles what has changed, and it leaves you in a Lisp
where you can mix Lisp and Dylan:

\begin{repl}
$ sbcl
* (load "src/load-dylan.lisp")

Call DYLAN:REP to start a Dylan read-eval-print loop.
;Dylan loaded.  (dylan:rep) starts a read-eval-print loop.
T
* (dylan:rep)

Entering Dylan read-eval-print-loop.
Leave by typing "(exit)"

?
\end{repl}

\noindent
There \code{dylan:done} returns to the Lisp prompt. What you defined is
not lost: if you call \code{(dylan:rep)} again, the variables, classes and
methods you defined earlier are still there. To leave SBCL altogether,
type \code{(exit)} at the \code{?} prompt, or \code{(sb-ext:exit)} at the
Lisp prompt.

\paragraph{When something goes wrong.} An error that your Dylan code does
not handle stops the loop and enters SBCL's debugger. The debugger
offers restarts that return you to the Dylan prompt; choose one by typing
its number:

\begin{repl}
? (error "Account ~a is overdrawn" 42)

debugger invoked on a DYLAN::DYLAN-CONDITION in thread
#<THREAD tid=701926 "main thread" RUNNING {1200030003}>:
  Error: Account 42 is overdrawn

Type HELP for debugger help, or (SB-EXT:EXIT) to exit from SBCL.

restarts (invokable by number or by possibly-abbreviated name):
  0: [DYLAN-RESTART] Return to the read-eval-print loop.
  1: [DYLAN-REP    ] Return to the Dylan read-eval-print loop.
  2: [ABORT        ] Return to the Dylan read-eval-print loop.
...
0] 0
;Aborted

? (+ 3 4)

;Value: 7
\end{repl}

\begin{notebox}
A few errors are detected by Common Lisp itself rather than by Dylan; for
example, \code{(car 5)} is reported as a Lisp \code{TYPE-ERROR}. For such
an error the debugger still offers \code{DYLAN-REP}, ``Return to the Dylan
read-eval-print loop'', which gives you a new \code{?} prompt. Dylan's own
handlers see these errors as instances of \code{<error>}, so
\code{handler-case} catches them normally. In a Lisp you loaded the
system into, the debugger's \code{ABORT} leaves the loop for the Lisp
\code{*} prompt; nothing you defined is lost; type \code{(dylan:rep)} to
carry on.
\end{notebox}

\paragraph{Running a file at the loop.} The loop reads from standard
input, so you can feed it a whole file of Dylan from the shell. This runs
every form in \file{program.dyl} and then leaves:

\begin{shell}
$ bin/dylan < program.dyl
\end{shell}

Each form's value is printed after \code{;Value:}, as if you had typed it.
An unhandled error prints the debugger's report and ends the run, since
the debugger finds no more input.

\section{Compiling a file to Lisp}\label{sec:start-compile-file}

\index{dylan:compile-file}
The interactive implementation can also compile a whole file of Dylan
into a file of Common Lisp, which can then be loaded and run without
reading or compiling any Dylan. Suppose \file{greet.dyl} contains this:

\begin{dylan}
(define-method greet ((name <string>))
  (display "Hello, ")
  (display name)
  (display "!")
  (newline))

(greet "world")
\end{dylan}

Load \file{src/load-compiler.lisp} and call \code{dylan:compile-file}
with the names of the input and output files:

\begin{shell}
$ sbcl --non-interactive --eval '(load "src/load-compiler.lisp")' \
       --eval '(dylan:compile-file "greet.dyl" "greet.lisp")'
\end{shell}

Compiling does not run the program; it only writes \file{greet.lisp}.
That file is ordinary Common Lisp. To run it, load
\file{src/load-runtime.lisp} and then load the compiled file:

\begin{shell}
$ sbcl --non-interactive --eval '(load "src/load-runtime.lisp")' \
       --eval '(load "greet.lisp")'
...
Hello, world!
\end{shell}

You can also load a compiled file into a session and then start the
loop; the file's definitions are then available at the \code{?} prompt,
and you can call them, redefine them, or add methods to them.

\begin{notebox}
In this version, \file{load-compiler.lisp} and \file{load-runtime.lisp}
both load the whole system, exactly as \file{load-dylan.lisp} does; they
are separate entry points so that a program can say what it needs.
Chapter~\ref{ch:lisp} describes \code{dylan:compile-file},
\code{dylan:compile-to-file} and modules in detail.
\end{notebox}

\section{A first program with dylanc}\label{sec:start-dylanc}

\index{dylanc!first program}
The compiler to C takes one or more files of Dylan and produces a native
executable. Write this into \file{hello.dyl}:

\begin{dylan}
(define-method greet ((name <string>))
  (format #t "Hello, ~a!~%" name))

(greet "world")
\end{dylan}

Then compile it and run the result:

\begin{shell}
$ bin/dylanc hello.dyl
$ ./hello
Hello, world!
\end{shell}

A program is simply a sequence of top-level forms, run in order from the
top of the file when the executable starts. There is no \code{main}
function; here the last form calls \code{greet}.

\begin{notebox}
The 1992 manual uses \code{format} in examples but does not define it;
both implementations provide it (Chapter~\ref{ch:printing}).
\end{notebox}

\paragraph{The first build takes longer.} The first time you run
\cmd{dylanc}, it compiles the compiler itself into fasl files, builds the
C runtime library with \cmd{make}, and compiles the Dylan library's C
code, all inside the clone (under \file{compiler/fasl/} and
\file{runtime/build/}). Each of these is kept and reused, so later
builds of a small program take only a second or two. The clone therefore
needs to be writable by the user who runs \cmd{dylanc}.

\paragraph{The executable's name.} By default the executable is named
after the first file, without the \file{.dyl}. The \cmd{-o} option names it
yourself:

\begin{shell}
$ bin/dylanc -o greeter hello.dyl
$ ./greeter
Hello, world!
\end{shell}

\paragraph{Several files.} A program can be split across several files.
Name them all on the command line; they are compiled, in the order given,
into one program, exactly as if they were one long file. Put definitions
before the code that uses them at top level. With \file{shapes.dyl}

\begin{dylan}
(define-class <square> (<object>)
  (side init-keyword: side:))

(define-method area ((s <square>))
  (* (side s) (side s)))
\end{dylan}

\noindent and \file{main.dyl}

\begin{dylan}
(print (area (make <square> side: 7)))
\end{dylan}

\noindent the program is built and run like this:

\begin{shell}
$ bin/dylanc -o app shapes.dyl main.dyl
$ ./app
49
\end{shell}

\paragraph{What files appear.} Normally the only file \cmd{dylanc} leaves
behind is the executable. The C it generates is written beside the
executable and deleted once the C compiler has finished with it. The \cmd{-k} (\cmd{-{}-keep-c}) option keeps it:

\begin{shell}
$ bin/dylanc -k hello.dyl
$ ls
hello  hello.c  hello.dyl  hello-dylan.c  hello-dylan-user.c  hello.h
\end{shell}

\file{hello.c} holds the program's data, its generic-function dispatch
and the C \code{main} function; \file{hello.h} is a header the other files
share; and there is one file for each module's functions:
\file{hello-dylan-user.c} for your code, which is in the module
\code{dylan-user}, and \file{hello-dylan.c} for the library's module
\code{dylan}. The generated C contains \code{\#line} directives that point
back to your Dylan source, so a C debugger or a C compiler's message
names the line of Dylan it came from. Reading this C is a good way to see
what the compiler did with your program; Chapter~\ref{ch:internals}
explains it.

\paragraph{Errors when compiling.} A syntax error stops the compilation
with a message giving the file, line and column, and a non-zero exit
status:

\begin{shell}
$ bin/dylanc bad.dyl
bad.dyl:1:15: error: End of file inside a Dylan list.
\end{shell}

A name that is used but never defined is only a warning; the program is
still built, and the error comes
when the program runs and reaches that name:

\begin{shell}
$ bin/dylanc -o bad main.dyl
main.dyl:1:8: warning: area is not defined (module dylan-user)
main.dyl:1:14: warning: <square> is not defined (module dylan-user)
$ ./bad
Dylan error: the variable area@dylan-user has no value
\end{shell}

Take these warnings seriously: they almost always mean a misspelling or a
missing file on the command line. A program that ends with an unhandled
error prints \code{Dylan error:} and a message to standard error and exits
with status~1.

\paragraph{Other options.} \cmd{bin/dylanc -h} lists every option. The
ones you are most likely to want soon are \cmd{-O0} to \cmd{-O3} (how hard
the C compiler optimizes; the default is \cmd{-O2}), \cmd{-Os} for a small
executable, \cmd{-g} for a debugging build that checks every inline
operation, and \cmd{-{}-trampoline} and \cmd{-{}-cc} for other C
compilers. Chapter~\ref{ch:dylanc} describes them all.

\section{The examples and the test suites}\label{sec:start-tests}

The repository includes the examples from the 1992 manual and a large
set of tests. You do not need to run any of them to use the system, but
they are a useful check that an installation works, and the examples are
good reading.

The interactive implementation runs the manual's examples, in
\file{dylan-examples.dyl}, and compares the results with the transcripts
in \file{src/reference/}. It must be run from the \file{src} directory:

\begin{shell}
$ cd src
$ sbcl --eval '(load "load-dylan.lisp")' \
       --eval '(load "full-test.lisp")' \
       --eval '(dylan:test-examples)'
\end{shell}

The compiler to C has its own suites, among them the manual's examples
chapter by chapter, tests of the code generator, stress tests, tests of
the \cmd{dylanc} command, and tests of the C runtime:

\begin{shell}
$ sbcl --script tests/conformance/run.lisp
$ sbcl --script tests/codegen/run.lisp
$ sbcl --script tests/stress/run.lisp
$ sh tests/driver/run.sh
$ make -C runtime test
\end{shell}

\begin{warnbox}
By default each \file{run.lisp} builds every test with both GCC and Clang,
with \code{musttail} and with trampolines, which takes a long time. Give
the mode \cmd{quick} as an argument, as in
\cmd{sbcl -{}-script tests/conformance/run.lisp quick}, for a faster run
with one configuration. Appendix~\ref{app:testing} describes the modes and
the benchmarks.
\end{warnbox}

\section{Where to go next}\label{sec:start-next}

Chapter~\ref{ch:tutorial} is a hands-on tour of the language. Start the
interactive loop, keep this book open beside it, and type the examples as
you read; at the end of the tutorial you will build a small program and
compile it with \cmd{dylanc}.

If you already know Lisp or Scheme and want the details first, the
language reference begins with Chapter~\ref{ch:syntax}. The interactive
implementation is described fully in Chapter~\ref{ch:lisp}, and the
compiler in Chapter~\ref{ch:dylanc}.
